\documentclass[11pt]{article}
\usepackage{estudio}
\cuadernillo{4 --- Resultados y conclusiones}

\begin{document}

\portada{4}{Resultados y conclusiones}{%
  Diapositivas 11 a 28 \textbullet\ expositores A, B y C \textbullet\ 7:50\\[4pt]
  Los catorce gráficos, por qué dieron así, y el cierre}

\tableofcontents
\clearpage

% ===========================================================================
\section{Cómo está armada la sección}

Resultados tiene \textbf{14 diapositivas de contenido} sobre 21 totales: es
dos tercios de la presentación y donde se juega la nota. El orden sigue el punto
2.4 de la guía de la cátedra: para cada modelo, primero animación, después
observable contra tiempo, después input contra observable.

\begin{center}
\begin{tabular}{clll}
\toprule
Exp. & Diapositivas & Contenido & Tiempo \\
\midrule
A & 11--16 & Vicsek: animación, $v_a(t)$, $v_a(\eta)$, $S(t)$, $S(\eta)$ & 2:40 \\
B & 17--21 & Votante: los mismos cinco & 2:30 \\
C & 22--28 & Comparaciones, $v_a$ vs $S$, benchmark, conclusiones & 2:40 \\
\bottomrule
\end{tabular}
\end{center}

Correspondencia con los incisos del enunciado: (a) diapositivas 12 y 17;
(b) 13 y 18; (c) 14, 19 y 22; (d) 15, 16, 20, 21 y 23; (e) 24; (f) se resuelve
superponiendo los dos modelos en 22, 23 y 24; (g) 25.

\ojo{\textbf{Promedio real: 34 segundos por diapositiva de gráfico.} Eso alcanza
para: decir qué hay en los ejes, señalar la tendencia, y dar el número clave. No
alcanza para describir la figura entera. Elegí \emph{una} cosa por diapositiva y
decila bien.}

% ===========================================================================
\section{Guion A: Vicsek (diapositivas 11 a 16)}

\diapo{11}{Divisoria: Resultados}{2 s}{A}

\diapo{12}{Vicsek: dinámica característica}{45 s}{A}

Dos animaciones lado a lado: $\eta = 1{,}0$ (ordenado) y $\eta = 4{,}0$
(desordenado), ambas a $\rho = 2$ ($N = 200$). Color según el ángulo de la
velocidad.

\begin{quote}
\itshape
``Estas son dos animaciones características del modelo de Vicsek, las dos a
densidad 2, o sea 200 partículas. Cada flecha es la velocidad de una partícula,
con origen en su posición, y el color codifica el ángulo. A la izquierda, con
ruido 1, se ve la fase ordenada: casi todas las flechas apuntan igual y el color
es prácticamente uniforme, la bandada se mueve como un bloque. A la derecha, con
ruido 4, el color está completamente mezclado: cada partícula va para donde
quiere y no hay movimiento colectivo. Entre esos dos extremos está la transición
que vamos a caracterizar.''
\end{quote}

\clave{El recurso visual más potente que tenés es el \textbf{color}: uniforme
= ordenado, mezclado = desordenado. Usalo, porque se entiende de inmediato y te
ahorra veinte segundos de explicación.}

\diapo{13}{Vicsek: evolución temporal de la polarización}{33 s}{A}

$v_a(t)$ a $\rho = 2$, tres niveles de ruido superpuestos, con líneas verticales
punteadas marcando el inicio de la ventana de medición.

\begin{quote}
\itshape
``Acá está la polarización en función del tiempo, a densidad 2, para tres
niveles de ruido. Las tres arrancan del mismo lugar, alrededor de 0,07, que es
el valor esperado para 200 partículas con direcciones al azar. Con ruido bajo la
polarización crece rápido y se estabiliza cerca de 1 en menos de 200 pasos; con
ruido alto se queda abajo, porque el sistema ya nace desordenado y no hay nada
que relajar. Las líneas verticales punteadas marcan dónde empieza la ventana
sobre la que promediamos, y son las que devuelve el criterio automático que
explicó [C]: el escalar que reportamos de acá en adelante es el promedio sobre
esa ventana.''
\end{quote}

\clave{Esta diapositiva justifica \emph{todas} las que siguen. Es donde
explicás de dónde sale el escalar. No la apures.}

\diapo{14}{Vicsek: polarización en función del ruido}{33 s}{A}

$v_a(\eta)$ para las densidades, con barras de error.

\begin{quote}
\itshape
``Esta es la curva principal del modelo de Vicsek: polarización en función del
ruido, para cada densidad, con las barras de error sobre 10 realizaciones. Las
tres curvas arrancan en 1 y decaen de forma monótona hacia el piso de tamaño
finito. Lo importante es que \emph{la curva se corre hacia la derecha al
aumentar la densidad}: el sistema tolera más ruido cuanto más denso es. Tomando
como referencia el cruce por 0,5, pasa de ruido 2,9 en densidad 2, a 3,2 en
densidad 4, a 3,4 en densidad 8. La razón es directa: con más vecinos, el
promedio circular atenúa mejor el ruido individual.''
\end{quote}

\dato{2,90 --- 3,19 --- 3,44 para $\rho = 2, 4, 8$. Son números medidos, no
redondeos: memorizalos.}

\diapo{15}{Vicsek: evolución temporal de la componente gigante}{22 s}{A}

$S(t)$ a $\rho = 1/\pi$ ($N = 32$).

\begin{quote}
\itshape
``Lo mismo para el segundo observable, la fracción de la componente gigante. Lo
mostramos a densidad $1/\pi$, o sea 32 partículas, y ya explico por qué: en las
densidades del enunciado esta curva es una línea recta en 1 y no muestra nada.
Acá sí se ve dinámica: con ruido bajo las partículas se agrupan y la componente
gigante crece; con ruido alto la red se fragmenta y $S$ cae.''
\end{quote}

\ojo{Anticipá vos el ``¿y las tres densidades del enunciado?'', porque si no te
lo van a preguntar. La respuesta está una diapositiva más adelante y en la nota
al costado de la 23: a $\rho = 2, 4, 8$ vale $S \approx 1$ para todo $\eta$.}

\diapo{16}{Vicsek: componente gigante en función del ruido}{25 s}{A}

$S(\eta)$.

\begin{quote}
\itshape
``Y esta es la curva de la componente gigante contra el ruido. En las
densidades del enunciado $S$ se queda arriba de 0,9 en todo el rango: el
observable está degenerado, no distingue la fase ordenada de la desordenada. La
razón es geométrica: el número medio de vecinos vale 6,3, 12,6 y 25,1, todos muy
por encima del umbral de percolación, que está en 4,5. La red percola sin
importar cuán ordenado esté el sistema. En las densidades subcríticas, en
cambio, sí discrimina: $S$ cae de casi 1 a menos de 0,2.''
\end{quote}

\clave{El número \textbf{4,5} es el umbral de percolación, y \textbf{6,3 /
12,6 / 25,1} son los $\langle k \rangle$ del enunciado. Con esos cuatro números
convertís ``el observable no dio nada'' en ``el observable dio exactamente lo
que la geometría predice''. Es la diferencia entre un resultado nulo y un
resultado explicado.}

Cerrá pasándole la palabra a B.

% ===========================================================================
\section{Guion B: modelo de votante (diapositivas 17 a 21)}

Tu bloque es paralelo al de A, así que \textbf{no repitas las definiciones}: los
ejes, el criterio de estacionario y las barras de error ya los explicaron. Vos
mostrás las diferencias.

\diapo{17}{Votante: dinámica característica}{40 s}{B}

Dos animaciones a $\rho = 2$: $\eta = 0{,}05$ (ordenado) y $\eta = 1{,}0$
(desordenado).

\begin{quote}
\itshape
``Las mismas dos situaciones para el modelo de votante, también a densidad 2.
Lo primero que hay que notar son los valores de ruido: para tener una bandada
ordenada acá hace falta ruido 0,05, mientras que en Vicsek alcanzaba con 1. Y ya
con ruido 1 el votante está completamente desordenado, cuando Vicsek a ese
mismo ruido seguía ordenado. La región ordenada del votante es casi un orden de
magnitud más angosta. Además, mirando el panel ordenado, se ve que el color
varía más entre partículas próximas que en Vicsek: copiar a un solo vecino no
promedia el ruido.''
\end{quote}

\clave{Tu bloque entero tiene una sola tesis: \emph{el votante pierde el orden
muchísimo antes}. Plantala acá, en la primera diapositiva, con los números
$0{,}05$ contra $1{,}0$.}

\diapo{18}{Votante: evolución temporal de la polarización}{28 s}{B}

\begin{quote}
\itshape
``La evolución temporal de la polarización del votante, a la misma densidad 2.
Comparada con la de Vicsek hay dos diferencias claras: converge mucho más lento,
y es mucho más ruidosa. Vicsek llegaba a polarización 0,8 en menos de 200 pasos;
el votante puede tardar más de mil, y la dispersión entre semillas es grande.
Eso es porque el orden no se construye promediando, se construye por
engrosamiento: dominios de orientación parecida que crecen copiándose entre
vecinos.''
\end{quote}

\diapo{19}{Votante: polarización en función del ruido}{32 s}{B}

\begin{quote}
\itshape
``La curva de polarización contra ruido del votante. La forma es la misma
---decae monótonamente--- pero la escala del eje horizontal es completamente
distinta: el colapso ocurre entre 0 y 0,5, no entre 2 y 4. Y hay algo que a
primera vista sorprende: acá la dependencia con la densidad está
\emph{invertida}. El cruce por 0,5 pasa de ruido 0,37 en densidad 2, a 0,29 en
densidad 4, a 0,19 en densidad 8. O sea que cuanto más denso, \emph{menos} ruido
tolera, exactamente al revés que Vicsek.''
\end{quote}

\ojo{\textbf{Esta es la diapositiva donde más probablemente te pregunten}, y la
pregunta va a ser ``¿por qué se invierte?''. La respuesta completa está en la
sección \ref{sec:inversion}. Estudiala hasta poder decirla sin pensar. En una
frase: \emph{el votante ordena por engrosamiento, y el tiempo de consenso crece
con $N$; a 2000 pasos, el sistema de 800 partículas todavía no terminó de
ordenarse ni siquiera con ruido cero}.}

\diapo{20}{Votante: evolución temporal de la componente gigante}{22 s}{B}

\begin{quote}
\itshape
``La componente gigante contra el tiempo para el votante, otra vez a densidad
$1/\pi$ por el mismo motivo que antes. El comportamiento es cualitativamente
igual al de Vicsek: se agrupa a ruido bajo y se fragmenta a ruido alto.''
\end{quote}

\diapo{21}{Votante: componente gigante en función del ruido}{28 s}{B}

\begin{quote}
\itshape
``Y la componente gigante contra el ruido. La conclusión es la misma que en
Vicsek, y esto es importante: en las densidades del enunciado $S$ se queda
arriba de 0,92 en todo el rango, degenerado, porque el argumento de percolación
es puramente geométrico y no depende de la regla de alineación. En las
densidades subcríticas vuelve a discriminar. O sea que el observable $S$ no
distingue entre los dos modelos: distingue entre regímenes de densidad.''
\end{quote}

\clave{Ese cierre ---``$S$ no distingue modelos, distingue densidades''--- es un
buen remate para tu bloque y le arma el pie a C, que arranca con las
comparaciones.}

% ===========================================================================
\section{Guion C: comparaciones y cierre (diapositivas 22 a 28)}

\diapo{22}{Comparación: polarización}{35 s}{C}

Los dos modelos superpuestos, solo las medias.

\begin{quote}
\itshape
``Acá están los dos modelos superpuestos en la misma figura, solo las medias
para que se lea. La diferencia de escala es lo que salta: las curvas de Vicsek
están corridas casi un orden de magnitud a la derecha respecto de las del
votante. Y las dos familias se comportan de manera opuesta con la densidad: en
Vicsek las curvas se corren a la derecha al aumentar la densidad, y en el
votante se corren a la izquierda. Ese cruce de tendencias es el resultado
central de la comparación.''
\end{quote}

\diapo{23}{Comparación: componente gigante}{30 s}{C}

Solo las densidades subcríticas, con la nota al costado sobre $\rho = 2, 4, 8$.

\begin{quote}
\itshape
``La comparación de la componente gigante, y acá mostramos solo las densidades
subcríticas por lo que ya se dijo: en las tres densidades del enunciado $S$ vale
aproximadamente 1 para todo ruido en los dos modelos, así que la figura sería
tres pares de líneas rectas. En el régimen subcrítico, en cambio, la caída es
clara, y los dos modelos dan curvas muy parecidas: la conectividad es un
fenómeno geométrico y no depende de cómo se alinean las partículas.''
\end{quote}

\diapo{24}{Polarización en función de la componente gigante}{33 s}{C}

Inciso (e). Cada punto es un valor de $\eta$.

\begin{quote}
\itshape
``El inciso (e) pide relacionar directamente los dos observables. Cada punto es
un valor de ruido. Y el resultado es negativo para las densidades del enunciado,
y hay que decirlo como tal: los puntos se apilan sobre la franja de $S$ cercano
a 1 y se extienden sobre todo el rango de polarización. O sea, el sistema
recorre desde totalmente ordenado hasta totalmente desordenado sin que la
conectividad cambie: orden direccional y conectividad espacial están
desacoplados. En las densidades subcríticas, en cambio, los puntos se alinean
sobre una diagonal: ahí sí, más orden va con más conectividad, porque las dos
cosas son manifestaciones del mismo agrupamiento dinámico.''
\end{quote}

\clave{Un resultado negativo bien explicado vale igual que uno positivo. Lo que
la cátedra castiga no es que $S$ no correlacione con $v_a$, es que no sepas
\emph{por qué}. Decilo con la palabra ``desacoplados'' y con el argumento de
percolación.}

\diapo{25}{Tiempo de ejecución del Cell Index Method}{30 s}{C}

Inciso (g), escala doble logarítmica.

\begin{quote}
\itshape
``El inciso (g) pide comparar los tiempos del CIM contra los del TP1. Medimos
exactamente la misma operación de los dos lados: la búsqueda de vecinos.
En TP2 está cronometrada dentro del paso temporal, alrededor de la
reconstrucción de la grilla, y no incluye ni el arranque del proceso ni la
escritura de la trayectoria. Igualamos las condiciones: misma caja, mismo radio,
mismo número de celdas y partículas puntuales en los dos binarios. El resultado
es que las dos series tienen la misma pendiente logarítmica, 1,43 y 1,46, o sea
el mismo escalado con el número de partículas. En valor absoluto el CIM de TP2
es más rápido, porque la grilla es un objeto persistente que no reasigna memoria
en cada llamada. Y la pendiente da superlineal en los dos casos porque el
barrido aumenta el número de partículas con la caja fija, así que también
aumenta la densidad.''
\end{quote}

\ojo{Ese último punto es el que te pueden pinchar: ``el CIM es $\mathcal{O}(N)$,
¿por qué la pendiente es 1,4?''. La respuesta es que la linealidad vale
\emph{a densidad constante}, y este barrido no la mantiene. Ver P16.}

\diapo{27}{Conclusiones}{25 s}{C}

Cinco puntos. \textbf{No los leas}: están escritos. Enunciá los tres que
importan y dejá que la cátedra lea el resto.

\begin{quote}
\itshape
``Para cerrar. En los dos modelos hay una transición de orden a desorden
controlada por el ruido. En Vicsek el ruido tolerado crece con la densidad,
porque promediar sobre más vecinos atenúa mejor el ruido; en el votante la
dependencia se invierte y la región ordenada es un orden de magnitud más
angosta, porque copiar a un solo vecino transmite el ruido íntegro. La
componente gigante resultó degenerada en las tres densidades del enunciado, por
un argumento de percolación puramente geométrico, y solo discrimina en el
régimen subcrítico. Y el Cell Index Method conserva el mismo escalado al
integrarse dentro del motor dinámico. Muchas gracias.''
\end{quote}

\ojo{\textbf{Verificar antes del viernes:} la diapositiva 27 cita para el
votante los valores $\eta = 0{,}34$, $0{,}26$ y $0{,}21$. Recalculados sobre
\texttt{data/sweep/summary.csv} dan \textbf{0,37, 0,29 y 0,19}. Los de Vicsek
(2,9 / 3,2 / 3,4) sí coinciden exactamente. La conclusión cualitativa ---la
inversión--- no cambia, pero los tres números del votante hay que corregirlos o
confirmar de dónde salieron.}

\diapo{28}{Muchas gracias}{5 s}{C}

Sin la palabra ``¿preguntas?'' (guía 2.6). Quedate en esta diapositiva mientras
responden.

% ===========================================================================
\section{Por qué nos dio eso}

Esta sección es el corazón del cuadernillo. Cada resultado con su mecanismo.

\subsection{Vicsek: por qué el ruido crítico crece con la densidad}

Es el resultado ``esperado'' y el argumento es cuantitativo. La partícula toma
el promedio circular de $k = \langle k \rangle$ direcciones, cada una con su
propio ruido independiente. Promediar $k$ variables aleatorias independientes
reduce la dispersión del promedio por un factor $k^{-1/2}$.

Entonces: más densidad $\Rightarrow$ más vecinos $\Rightarrow$ mejor filtrado
del ruido $\Rightarrow$ hace falta más $\eta$ para destruir el orden. Los datos:

\begin{center}
\begin{tabular}{lcccccc}
\toprule
$N$ & 11 & 16 & 32 & 200 & 400 & 800 \\
$\langle k \rangle$ & 0{,}33 & 0{,}50 & 1{,}00 & 6{,}3 & 12{,}6 & 25{,}1 \\
\midrule
$\eta$ con $v_a = 0{,}5$ & 2{,}07 & 2{,}07 & 2{,}22 & 2{,}90 & 3{,}19 & 3{,}44 \\
\bottomrule
\end{tabular}
\end{center}

\dato{Monótono creciente sobre las seis densidades, sin una sola excepción. Es
el comportamiento de libro.}

\subsection{El votante y la inversión con la densidad}
\label{sec:inversion}

\textbf{Este es el resultado que hay que saber defender.} Los datos:

\begin{center}
\begin{tabular}{lcccccc}
\toprule
$N$ & 11 & 16 & 32 & 200 & 400 & 800 \\
\midrule
$\eta$ con $v_a = 0{,}5$ & 1{,}03 & 0{,}79 & 0{,}70 & 0{,}37 & 0{,}29 & 0{,}19 \\
\bottomrule
\end{tabular}
\end{center}

Monótono \emph{decreciente} sobre las seis densidades: exactamente al revés que
Vicsek. ¿Por qué?

\textbf{Primero, por qué el votante no gana nada con más vecinos.} En Vicsek, el
promedio sobre $k$ vecinos atenúa el ruido. En el votante no hay promedio: la
partícula copia a \emph{uno} y hereda su ángulo tal cual, con el ruido íntegro.
Aumentar $k$ no mejora el filtrado, porque no hay filtrado. Eso ya explica por
qué el votante es mucho más frágil que Vicsek a cualquier densidad.

\textbf{Segundo, por qué empeora con la densidad.} Acá está la clave, y es un
efecto de \emph{tiempo finito}. El votante no ordena por promediado sino por
\emph{engrosamiento}: se forman dominios de orientación parecida que crecen
copiándose entre vecinos, como en un proceso de coalescencia de dominios. El
tiempo que tarda el sistema en llegar al consenso \textbf{crece con el número de
partículas}. Como todas las corridas son de 2000 pasos, el sistema de 800
partículas está mucho más lejos del consenso que el de 200 al momento de medir.

\dato{\textbf{La prueba, medida sobre nuestros propios datos y con ruido
exactamente cero} ($\eta = 0$, o sea sin nada que desordene):
\vspace{4pt}

\begin{tabular}{lcccc}
\toprule
$N$ & 11 & 16 & 32 & 800 \\
\midrule
$v_a$ en $t = 2000$ & 1{,}000 & 0{,}989 & 0{,}999 & \textbf{0{,}978 $\pm$ 0{,}052} \\
\bottomrule
\end{tabular}
\vspace{6pt}

Con \emph{cero} ruido, el sistema de 800 partículas no llega al consenso en 2000
pasos, y con una dispersión entre semillas enorme. Los sistemas chicos sí. La
curva completa de $N=800$ a $\eta = 0$: $v_a$ vale $0{,}35$ en $t = 100$,
$0{,}62$ en $t = 500$, $0{,}85$ en $t = 1000$ y $0{,}98$ en $t = 2000$: todavía
está subiendo cuando termina la corrida.}

\clave{La respuesta de una frase: \emph{``el votante ordena por engrosamiento,
no por promediado, y el tiempo de consenso crece con $N$; a 2000 pasos el
sistema de 800 partículas todavía no terminó de ordenarse ni siquiera con ruido
cero, así que su polarización medida es más baja y la curva parece corrida hacia
ruidos menores''}.}

\ojo{Cuidado con cómo lo enunciás. No digas ``el ruido crítico del votante baja
con la densidad'' como si fuera una propiedad del punto crítico: lo honesto es
decir que \emph{a tiempo de simulación fijo}, el sistema más grande está menos
polarizado, y que eso corre aparentemente la transición. Distinguir un efecto de
tiempo finito de una propiedad del límite termodinámico es exactamente el tipo
de rigor que la cátedra premia.}

\teoria{Por qué el tiempo de consenso crece con $N$: en el modelo de votante
clásico, el consenso se alcanza por fluctuaciones de dominio, y en dos
dimensiones el tiempo característico escala como $N \ln N$. El mecanismo es
difusivo ---las paredes entre dominios hacen una caminata aleatoria--- mientras
que en Vicsek el promediado alinea globalmente en un tiempo que casi no depende
de $N$. Nuestros datos son consistentes con eso, aunque no lo medimos con
suficientes valores de $N$ como para verificar el exponente, y por eso no lo
afirmamos en la presentación.}

Hay además una segunda diferencia dinámica, visible en las diapositivas 13 y 18:
la convergencia del votante es mucho más lenta y más dispersa. A $\rho = 8$ y
$\eta = 0{,}2$, Vicsek llega a $v_a \geq 0{,}8$ entre los pasos 10 y 174 en las
10 semillas y termina en $0{,}998 \pm 0{,}000$; el votante lo alcanza recién
entre los pasos 454 y 1596, solo en 7 de las 10 semillas, y termina en
$0{,}379 \pm 0{,}188$.

\subsection{Por qué $S \approx 1$ en las densidades del enunciado}

Es geometría pura, y no tiene nada que ver con la dinámica de alineación.

Para partículas distribuidas al azar con densidad $\rho$ y radio de conexión
$r_c$, el número medio de vecinos es $\langle k \rangle = \rho \pi r_c^2$. Ese
es exactamente el parámetro de control de la \emph{percolación del grafo
geométrico aleatorio} en dos dimensiones, cuyo umbral está en
$\langle k \rangle_c \approx 4{,}5$.

\begin{center}
\begin{tabular}{lccc|ccc}
\toprule
$\rho$ & 2 & 4 & 8 & $1/\pi$ & $1/2\pi$ & $1/3\pi$ \\
$\langle k \rangle$ & 6{,}3 & 12{,}6 & 25{,}1 & 1{,}00 & 0{,}50 & 0{,}33 \\
\midrule
$S$ mínimo (Vicsek) & 0{,}91 & 0{,}99 & 1{,}00 & 0{,}15 & 0{,}17 & 0{,}17 \\
$S$ mínimo (votante) & 0{,}92 & 1{,}00 & 1{,}00 & 0{,}16 & 0{,}17 & 0{,}17 \\
\bottomrule
\end{tabular}
\end{center}

Las tres del enunciado están \emph{todas} por encima del umbral, así que la red
percola incluso antes de considerar la dinámica: aunque las partículas
estuvieran distribuidas completamente al azar, ya formarían un cluster gigante.
Por eso $S$ no distingue fases ahí. Las tres subcríticas están holgadamente por
debajo, y ahí toda componente gigante que aparezca es producto del
\emph{agrupamiento dinámico}.

\clave{La forma correcta de presentar esto no es ``$S$ no nos dio nada'', es
``\emph{$S$ dio exactamente lo que predice la percolación geométrica, y por eso
agregamos densidades subcríticas donde el observable sí discrimina}''. Un
resultado degenerado que sabés predecir es un resultado.}

Detalle fino, por si preguntan por la forma de la curva: en $\rho = 2$, $S$ baja
levemente hasta un mínimo de $0{,}905$ alrededor de $\eta \approx 3{,}8$ y
después \emph{vuelve a subir} hasta $0{,}991$ a ruido máximo. No es monótona.
La interpretación es geométrica: a ruido máximo las partículas se distribuyen de
manera casi uniforme, y a densidad supercrítica una distribución uniforme
percola igual de bien que una bandada compacta. El mínimo está en el medio,
donde el sistema forma bandas densas y deja regiones vacías.

\subsection{Por qué $v_a$ y $S$ están desacoplados (inciso e)}

Es consecuencia directa de lo anterior. En las densidades del enunciado, $S$
está clavado cerca de 1 mientras $v_a$ recorre todo su rango: la nube de puntos
es una franja horizontal. La correlación lineal es despreciable y ni siquiera es
monótona, porque $S$ baja y vuelve a subir.

En las densidades subcríticas los puntos se ordenan sobre una diagonal, porque
ahí las dos cosas ---orden direccional y conectividad--- son manifestaciones del
mismo agrupamiento: a ruido bajo las partículas se alinean, se juntan y forman
un cluster; al aumentar el ruido la bandada se disgrega y las dos cantidades
caen juntas.

\subsection{El benchmark del CIM (inciso g)}

Qué se mide, exactamente lo mismo de los dos lados: la búsqueda de vecinos.

\begin{itemize}
  \item \textbf{TP2}: cronómetro adentro de \texttt{step()}, justo antes y justo
    después de la reconstrucción de la grilla. No entra el arranque del proceso,
    ni la generación de partículas, ni la escritura de la trayectoria, ni las
    reconstrucciones extra que hace el motor para medir $S$.
  \item \textbf{TP1}: la medición de la búsqueda pura que ya traía su binario.
\end{itemize}

Cómo se igualan las condiciones: se fuerza $L = 10$, $r_c = 1$, contorno
periódico y el mismo $M$ en los dos binarios, y TP1 corre con radio de partícula
cero para que el predicado de vecindad sea centro a centro igual que en TP2. Los
dos descartan el 1\% de las mediciones más lentas, que es la regla que ya traía
TP1.

Resultados:

\begin{center}
\begin{tabular}{lcc}
\toprule
 & TP1 & TP2 \\
\midrule
Pendiente log-log (todos los $N$) & 1{,}43 & 1{,}46 \\
Tiempo en $N = 200$ & 0{,}042 ms & 0{,}029 ms \\
Tiempo en $N = 800$ & 0{,}396 ms & 0{,}329 ms \\
\bottomrule
\end{tabular}
\end{center}

Dos lecturas. \textbf{Primera}: las pendientes son prácticamente idénticas, o
sea que integrar el CIM dentro de un motor dinámico no degrada el escalado.
\textbf{Segunda}: en valor absoluto TP2 es entre un 20\% y un 50\% más rápido,
porque la grilla es un objeto persistente que vacía y rellena sus buffers en
lugar de reasignarlos en cada llamada, como hacía TP1.

\ojo{\textbf{Por qué la pendiente no es 1.} El CIM es $\mathcal{O}(N)$
\emph{a densidad constante}. Este barrido aumenta $N$ de 10 a 1000 con la caja
fija en $L = 10$, así que la densidad va de $0{,}1$ a $10$ y el número medio de
vecinos por partícula pasa de $0{,}3$ a $31$. Hay más trabajo por partícula a
medida que crece $N$, y por eso las dos series salen superlineales. Es una
consecuencia del diseño del barrido, no un defecto del método.}

% ===========================================================================
\section{Preguntas y respuestas}

\subsection{Sobre la transición y la comparación de modelos}

\pregunta{¿Por qué en Vicsek el ruido crítico crece con la densidad?}
\respuesta{Porque la regla promedia. Promediar $k$ direcciones con ruido
independiente reduce la dispersión del promedio por un factor $k^{-1/2}$, así
que con más vecinos el ruido individual se filtra mejor y hace falta más $\eta$
para destruir el orden. Medido en el cruce por $v_a = 0{,}5$: 2,90 para
$\rho=2$, 3,19 para $\rho=4$ y 3,44 para $\rho=8$, y sigue siendo monótono si se
incluyen las tres densidades subcríticas.}

\pregunta{¿Por qué en el votante la dependencia con la densidad está invertida?}
\respuesta{Por dos motivos que se suman. Primero, el votante no promedia, así
que más vecinos no le dan ninguna ventaja: copia a uno solo y hereda el ruido
íntegro. Y segundo, que es lo que explica la inversión: el votante ordena por
engrosamiento de dominios, y el tiempo hasta el consenso crece con el número de
partículas. Como todas las corridas son de 2000 pasos, el sistema de 800
partículas está mucho más lejos del consenso que el de 200 al momento de medir.
Lo verificamos con ruido exactamente cero: a $\eta = 0$ y $t = 2000$, los
sistemas de 11, 16 y 32 partículas llegan a $v_a \approx 1$, mientras que el de
800 se queda en $0{,}978$ con una dispersión de $0{,}05$, y su curva todavía
está subiendo cuando termina la corrida.}

\pregunta{Entonces, ¿el ruido crítico del votante realmente baja con la
densidad, o es un artefacto?}
\respuesta{Lo que medimos es que \emph{a tiempo de simulación fijo} el sistema
más grande está menos polarizado, y eso corre la curva hacia ruidos menores. No
afirmamos que sea una propiedad del punto crítico en el límite termodinámico:
para eso habría que correr hasta el consenso en cada tamaño, o hacer un análisis
de escaleo de tamaño finito, que no está en el alcance del TP. Lo reportamos
como lo que es, un efecto de tiempo finito, y lo respaldamos con la medición a
ruido cero.}

\pregunta{¿Por qué el votante es tanto más frágil al ruido que Vicsek?}
\respuesta{Porque no atenúa el ruido. Vicsek promedia sobre todo el vecindario y
el ruido se cancela parcialmente; el votante copia a un vecino y se lleva su
ruido completo. Por eso la región ordenada del votante es casi un orden de
magnitud más angosta: a densidad 2, Vicsek cruza $v_a = 0{,}5$ en $\eta = 2{,}9$
y el votante en $\eta = 0{,}37$.}

\pregunta{¿Por qué la convergencia del votante es más lenta y más dispersa?}
\respuesta{Porque el mecanismo es distinto. Vicsek alinea globalmente en cada
paso vía el promedio; el votante forma dominios locales que crecen copiándose,
y las paredes entre dominios se mueven de forma difusiva. Medido a $\rho = 8$ y
$\eta = 0{,}2$: Vicsek llega a $v_a \geq 0{,}8$ entre los pasos 10 y 174 en las
10 semillas y termina en $0{,}998$; el votante recién entre los pasos 454 y
1596, y solo en 7 de las 10 semillas.}

\pregunta{Si corrieran más pasos, ¿cambiaría el resultado del votante?}
\respuesta{Sí, y es la consecuencia honesta de nuestro propio análisis: con
corridas más largas el sistema de 800 partículas alcanzaría el consenso y su
curva se correría hacia ruidos mayores, probablemente hasta recuperar la
tendencia creciente con la densidad. No lo corrimos porque el barrido completo
son 4680 simulaciones, pero está identificado como el paso siguiente.}

\subsection{Sobre la componente gigante}

\pregunta{¿Por qué $S$ vale casi 1 en todas las densidades del enunciado?}
\respuesta{Por percolación geométrica. El número medio de vecinos es
$\langle k \rangle = \rho\pi r_c^2$, y ese es el parámetro de control de la
percolación del grafo geométrico aleatorio en 2D, cuyo umbral está en
$\approx 4{,}5$. Las tres densidades del enunciado dan 6,3, 12,6 y 25,1: todas
por encima. La red percola aunque las partículas estén distribuidas
completamente al azar, así que $S$ no puede distinguir la fase ordenada de la
desordenada.}

\pregunta{Entonces $S$ no sirve como observable.}
\respuesta{No sirve \emph{en ese régimen}, y saber por qué es el resultado. Por
eso agregamos tres densidades subcríticas, con $\langle k \rangle$ de 1, 0,5 y
0,33, holgadamente por debajo del umbral: ahí una configuración uniforme no
percola, así que toda componente gigante que aparezca es producto del
agrupamiento dinámico, y el observable sí discrimina entre fases: $S$ cae de
casi 1 a menos de 0,2.}

\pregunta{¿Por qué mostraron $S(t)$ solo a $\rho = 1/\pi$ y no a las tres
densidades del enunciado?}
\respuesta{Porque a $\rho = 2, 4, 8$ la curva es una línea recta en 1: $S$ ya
vale prácticamente 1 en la condición inicial y se queda ahí toda la corrida, no
hay transitorio de conectividad que mostrar. La figura no aportaría información
y por eso elegimos la densidad donde sí hay dinámica. El valor
constante en las densidades del enunciado está indicado explícitamente al
costado de la diapositiva 23, y las seis series están en el informe.}

\pregunta{La curva de $S(\eta)$ a $\rho = 2$ no es monótona, ¿por qué?}
\respuesta{Baja hasta un mínimo de $0{,}905$ cerca de $\eta \approx 3{,}8$ y
después vuelve a subir hasta $0{,}991$. La interpretación es geométrica: a ruido
máximo las partículas quedan distribuidas de manera casi uniforme, y a densidad
supercrítica una distribución uniforme percola igual de bien que una bandada
compacta. El mínimo está en el medio, donde el sistema forma bandas densas y
deja regiones vacías, que es lo que más fragmenta la red.}

\pregunta{¿$S$ distingue entre los dos modelos?}
\respuesta{No, y es esperable: la conectividad es un fenómeno geométrico que
depende de dónde están las partículas, no de cómo deciden su dirección. Las
curvas de los dos modelos son muy parecidas en los dos regímenes. Lo que $S$
distingue es el régimen de densidad, no el modelo.}

\subsection{Sobre la relación $v_a$--$S$ (inciso e)}

\pregunta{¿Qué concluyen del gráfico de $v_a$ contra $S$?}
\respuesta{Que en las densidades del enunciado los dos observables están
desacoplados: el sistema recorre desde $v_a \approx 1$ hasta $v_a \approx 0$ sin
que $S$ se mueva de la franja cercana a 1. La correlación es despreciable y ni
siquiera es monótona. En las densidades subcríticas, en cambio, los puntos se
ordenan sobre una diagonal, porque ahí orden direccional y conectividad son
manifestaciones del mismo agrupamiento dinámico. Es un resultado negativo en un
régimen y positivo en el otro, y la causa de la diferencia es el umbral de
percolación.}

\subsection{Sobre el benchmark}

\pregunta{¿Qué midieron exactamente en el inciso (g)?}
\respuesta{La búsqueda de vecinos, la misma operación de los dos lados. En TP2
está cronometrada adentro del paso temporal, alrededor de la reconstrucción de
la grilla: no incluye el arranque del proceso, ni la generación de partículas,
ni la escritura de la trayectoria, ni las reconstrucciones adicionales que
hacemos para medir $S$. En TP1 es la medición de la búsqueda pura que ya traía
su binario.}

\pregunta{¿Cómo se aseguran de que la comparación sea justa?}
\respuesta{Forzamos los mismos parámetros en los dos binarios: $L = 10$,
$r_c = 1$, contorno periódico y el mismo $M$. Además corremos TP1 con radio de
partícula cero, para que su predicado de vecindad sea centro a centro igual que
el de TP2, y para que el $M$ máximo válido coincida. Y los dos descartan el 1\%
de las mediciones más lentas, que es la regla que ya usaba TP1.}

\pregunta{El CIM es $\mathcal{O}(N)$, ¿por qué la pendiente da 1,4 y no 1?}
\respuesta{Porque la linealidad vale a densidad constante, y este barrido no la
mantiene: aumentamos $N$ de 10 a 1000 con la caja fija en $L = 10$, así que la
densidad va de 0,1 a 10 y el número medio de vecinos por partícula pasa de 0,3 a
31. Con más partículas por celda, cada una hace más comparaciones, y el
escalado resulta superlineal. Lo relevante para el inciso es que las dos series
tienen la \emph{misma} pendiente ---1,43 y 1,46---, o sea que integrar el CIM
dentro del motor dinámico no degrada el escalado.}

\pregunta{¿Por qué el CIM de TP2 es más rápido que el de TP1?}
\respuesta{Porque la grilla es un objeto persistente: los vectores de celdas y
de listas de vecinos se vacían y se vuelven a llenar en cada paso, sin liberar
ni volver a reservar memoria. El TP1 creaba y destruía sus buffers en cada
llamada. Sobre el orden de $10^6$ reconstrucciones que hace el barrido, esa
diferencia se nota: entre un 20\% y un 50\% menos de tiempo.}

\subsection{Preguntas transversales}

\pregunta{¿Por qué $v_a$ no llega a cero a ruido máximo?}
\respuesta{Por tamaño finito. En la fase desordenada las direcciones son
independientes, así que la suma vectorial es una caminata aleatoria y
$v_a \approx N^{-1/2}$. Para $N = 200$ eso predice 0,071 y medimos 0,063; para
$N = 800$ predice 0,035 y medimos 0,031. El acuerdo es muy bueno en las seis
densidades.}

\pregunta{¿Determinaron el ruido crítico?}
\respuesta{Usamos el cruce por $v_a = 0{,}5$ como referencia operativa para
comparar las curvas, que es lo que mostramos. Calculamos además la
susceptibilidad y su máximo, pero no lo incluimos en la presentación porque con
estos tamaños el pico es ancho y chato y la posición del máximo no es estable
entre subconjuntos de semillas: reportarlo con tres cifras sugeriría una
precisión que los datos no respaldan. El enunciado tampoco lo pide.}

\pregunta{¿Cómo saben que las curvas no cambiarían con otra semilla?}
\respuesta{Cada punto es el promedio de 10 realizaciones independientes con
semillas derivadas por hash, y las barras de error son el desvío entre esas 10
medias. Fuera de la región de transición las barras son muy chicas; donde
crecen es cerca del punto crítico, que es exactamente donde se espera que el
sistema fluctúe más.}

\pregunta{¿Por qué las animaciones son solo a $\rho = 2$?}
\respuesta{Porque es la densidad donde la dinámica se ve mejor: con 200
partículas se distinguen las flechas individuales y a la vez hay suficiente
estructura colectiva. A $\rho = 8$ son 800 flechas y la figura se satura. El
enunciado pide animaciones para ``pocas situaciones características'', y elegimos
mostrar el contraste ordenado/desordenado, que es lo que hace legible la
transición.}

\pregunta{¿Qué harían distinto si tuvieran más tiempo?}
\respuesta{Tres cosas concretas. Correr el votante mucho más allá de 2000 pasos
para separar el efecto de tiempo finito del comportamiento asintótico, que es la
limitación principal que identificamos. Hacer un análisis de escaleo de tamaño
finito a densidad constante, variando $L$ y $N$ juntos, para poder hablar del
punto crítico en el límite termodinámico. Y correr el benchmark a densidad
constante, para verificar el escalado lineal del CIM en lugar de solo comparar
pendientes entre TP1 y TP2.}

\end{document}
